\begin{dang}{Lượng chất đã phân rã – lượng chất tạo thành}
\Anlythuyet{
\begin{itemize}
		\item Số hạt nhân đã phân rã của chất phóng xạ:
	\begin{align*}
		\boder{\Delta N=N_{0}-N=N_{0}\pbrace{1-e^{-\lambda t}}=N_{0}\pbrace{1-2^{-\tfrac{t}{T}}}}
	\end{align*}
	\item Khối lượng hạt nhân đã phân rã của chất phóng xạ:
	\begin{align*}
		\boder{\Delta m=m_{0}-m=m_{0}\pbrace{1-e^{-\lambda t}}=m_{0}\pbrace{1-2^{-\tfrac{t}{T}}}}
	\end{align*}
\item Tỉ số giữa lượng chất phân rã và lượng chất còn lại:
\begin{align*}
	\boder{\dfrac{\Delta N}{N}=\dfrac{\Delta m}{m}=2^{\tfrac{t}{T}}-1}
\end{align*}
	\item Số hạt nhân con - khối lượng hạt nhân con tạo thành: 
	\begin{align*}
		  \Hn{{Z_X}}{{A_X}}{X}\rightarrow \Hn{{Z_Y}}{{A_Y}}{Y}+\text{ tia phóng xạ}
	  \end{align*} 
	Ta thấy một hạt nhân mẹ phân rã thì sẽ có một hạt nhân con tạo thành nên
	\begin{itemize}
		\item Số hạt nhân con tạo thành chính là số hạt nhân mẹ đã phân rã: 
		$N_Y=\Delta N_X$
		\item Khối lượng hạt nhân con tạo thành: \boder{\bth{m_Y=\dfrac{N_Y.A_Y}{N_A}=\dfrac{\Delta N_X.A_Y}{N_A}=\dfrac{\Delta m_X.A_2}{A_1}}}
	\end{itemize}
\end{itemize}
}
\end{dang}
